\section{Introduction}

A visible cycle may close while a receipt fiber continues to evolve. In the
binary case, a nontrivial return receipt doubles the cycle length. The natural
next question is whether doubling is special to a two-element receipt or is
one instance of a general finite law.

This paper answers that question for an arbitrary finite receipt group $R$.
The local lift appears to contain one independent group element on every
visible edge. Gauge removes almost all of this apparent freedom. What remains
is a single group element up to conjugation: the ordered receipt accumulated
around the visible circuit.

The order of that holonomy element determines the entire lifted cycle profile.
If the receipt has order $d$, complete lifted return requires $d$ visible
circuits. Every lifted component therefore has length $nd$, and the number of
components is $|R|/d$. The argument is elementary and self-contained, although
it belongs to the standard setting of regular voltage-graph covers
\cite{gross-tucker}.

